\documentclass{beamer}
\usepackage{bm}
\usepackage{tabularx}
\setbeamertemplate{caption}[numbered]
\usepackage{xcolor}
\usepackage{multirow}
\usepackage{movie15}
\usepackage{Sweave}
%\usepackage{nameref}% http://ctan.org/pkg/nameref
%\newcommand{\currentchapter}{}
%\let\oldchapter\chapter
%\renewcommand{\chapter}[1]{\oldchapter{#1}\renewcommand{\currentchapter}{#1}}
\renewcommand{\thesection}{2.\arabic{section}}
\newcommand{\currentsection}{}
\let\oldsection\section
\renewcommand{\section}[1]{\oldsection{#1}\renewcommand{\currentsection}{#1}}

\newcommand{\currentsubsection}{}
\let\oldsubsection\subsection
\renewcommand{\subsection}[1]{\oldsubsection{#1}\renewcommand{\currentsubsection}{#1}}

\newcommand{\currentsubsubsection}{}
\let\oldsubsubsection\subsubsection
\renewcommand{\subsubsection}[1]{\oldsubsubsection{#1}\renewcommand{\currentsubsubsection}{#1}}

\newcommand{\boldbeta}{{\boldsymbol \beta}}
\newcommand{\boldtheta}{{\boldsymbol \theta}}
\newcommand{\bolddelta}{{\boldsymbol \delta}}
\newcommand{\bR}{\mathbb{R}}

\newenvironment{wideitemize}%
  {\begin{itemize}%
%    \setlength{\itemsep}{0pt}%
    \setlength{\parskip}{1em}}%
  {\end{itemize}}

%\usetheme{CambridgeUS}
\usetheme{Madrid}

%\title[Section~1.\insertsectionnumber \currentsection]{Chapter 1. Motivating Examples}
\title[\currentsection]{Chapter 2.Univariate Generalised Linear Models (GLM)}
\author[univariate GLM]{MAST90084 Statistical Modelling}
\vskip 1cm
\institute[Ch2]{Dennis Leung \\ \vskip 0.3cm SCHOOL OF MATHEMATICS AND STATISTICS\\
\vskip 0.3cm
   THE UNIVERSITY OF MELBOURNE}
\vskip 1cm
\date[Dennis Leung \hspace{8pt}]{}


\begin{document}

\frame{\titlepage}



\begin{frame}[fragile] \frametitle{Outline}
\tableofcontents
\end{frame}

\section{\S2.2 Likelihood Inference}\label{sec:2.2}
\begin{frame}{\S2.2 Likelihood Inference}
\begin{wideitemize}
\item
Given its full distributional assumption, regression analysis with GLMs can be based on the full \textbf{likelihood}. 
\item
Inference tasks include
\begin{itemize}
\item parameter point estimation
\item constructing confidence intervals
\item hypothesis testing
\item  goodness-of-fit tests 
\end{itemize}
\item
For now we assume the model is completely and correctly specified, i.e., the GLM assumptions are all ``true" (to be relaxed when we cover quasi-likelihood).  
\end{wideitemize}
\end{frame}

\subsection{\S2.2.1 Maximum likelihood estimation}\label{sec:2.2.1}
\begin{frame}{\S2.2.1 Maximum likelihood estimation (MLE)}
\begin{wideitemize}
\item
Data:  response $y_i$, covariates $\textbf{x}_i$ and design
vector $\textbf{z}_i = {\bf z}_i(\textbf{x}_i)$, $i=1,2,\cdots, n$.
%\item
%Model:  GLM where $E(y_i|\textbf{x}_i)=\mu_i=h(\eta_i)=h(\textbf{z}_i^T\!\mbox{\boldmath $\beta$})$,
%with response function $h(\cdot)$ and unknown coefficient $p\times 1$ vector $\mbox{\boldmath $\beta$}$. 
\item
Assumptions: 
\begin{wideitemize}
\item  design matrix $Z = Z_{n\times p}=(\textbf{z}_1,\cdots,\textbf{z}_n)^T$ has (full) rank $p$
\item  scale/dispersion parameter $\phi$ is
assumed known or otherwise consistently estimated. 
\end{wideitemize}

\item 
(Recall : we treat the design vectors as non-random)

\item We will estimate $\mbox{\boldmath $\beta$}$ by the {\bf maximum likelihood estimate} (MLE) $\hat{\boldsymbol \beta}$. 
\end{wideitemize}

\end{frame}

\begin{frame}{\S2.2.1 MLE (log-likelihood)}

%The focus is estimating $\mbox{\boldmath $\beta$}$. 
The MLE maximizes the log-likelihood, treated as a function in $\boldsymbol \beta$. Now the log-likelihood function is {\footnotesize
\begin{eqnarray*}
\ell(\mbox{\boldmath $\beta$})&=&\sum_{i=1}^n\ell_i(\mbox{\boldmath $\beta$})
=\sum_{i=1}^n \log f(y_i|\theta_i,\phi,\omega_i) \\
&=&\sum_{i=1}^n \frac{y_i\theta_i-b(\theta_i)}{\phi}\cdot\omega_i
+(\mbox{terms not involving $\theta_i$'s} )\\
&=& \sum_{i=1}^n \frac{y_i\theta (\mu_i)-b(\theta(\mu_i))}{\phi}\cdot\omega_i  + (ditto)\\
&=&\sum_{i=1}^n \frac{y_i\theta (h(\textbf{z}_i^T\!\mbox{\boldmath $\beta$}))-
b(\theta(h(\textbf{z}_i^T\!\mbox{\boldmath $\beta$})))}{\phi}\cdot\omega_i + (ditto)
\end{eqnarray*}}



 We find $\hat{\boldsymbol \beta}$ by solving for the root  of the derivative of $\ell$. 
\end{frame}

\begin{frame}{\S2.2.1  MLE (Score function)}
\begin{itemize}
\item
The \textbf{score function} is the first derivative of $\ell$: {\footnotesize
\begin{eqnarray*}
\textbf{s}(\mbox{\boldmath $\beta$})&\!\!\!=\!\!\!& \frac{\partial \ell(\mbox{\boldmath $\beta$})}{\partial\mbox{\boldmath 
$\beta$}}=\sum_{i=1}^n\frac{\partial \ell_i(\mbox{\boldmath $\beta$})}{\partial\mbox{\boldmath $\beta$}}
=\sum_{i=1}^n  \frac{\partial \ell_i(\mbox{\boldmath $\beta$})}{\partial\theta_i}\cdot
\frac{\partial \theta_i}{\partial\mu_i}\cdot \frac{\partial \mu_i}{\partial\eta_i} \cdot
\frac{\partial \eta_i}{\partial\mbox{\boldmath $\beta$}}\\
&\!\!\!=\!\!\!&\sum_{i=1}^n \left(\frac{y_i-b^\prime(\theta_i)}{\phi}\cdot\omega_i\right)\cdot
\left[v(h(\textbf{z}_i^T\!\mbox{\boldmath $\beta$}))\right]^{-1}\cdot
 \frac{\partial h(\eta_i)}{\partial\eta_i}\cdot \underbrace{ \textbf{z}_i }_{ = \frac{\partial \eta_i}{ \partial {\boldsymbol \beta}}}\\
&\!\!\!=\!\!\!& \sum_{i=1}^n \textbf{z}_i\cdot \frac{\partial h(\eta_i)}{\partial\eta_i}\cdot
\left[ \underbrace{\frac{\phi v(h(\textbf{z}_i^T\!\mbox{\boldmath $\beta$}))}{\omega_i}}_{\sigma^2_i({\boldsymbol \beta})}\right]^{-1}\cdot
\left[y_i-\mu_i(\mbox{\boldmath $\beta$})\right]
\end{eqnarray*}},
where
\begin{align*}
& \mu_i=b^\prime(\theta_i) \Rightarrow \quad \frac{\partial \mu_i}{\partial\theta_i}=
b^{\prime\prime}(\theta_i)=v(\mu_i)=v(h(\textbf{z}_i^T\!\mbox{\boldmath $\beta$})) \\
&\Rightarrow \quad   \frac{\partial \theta_i}{\partial\mu_i} 
\underset{\Big(\substack{inverse \\ function\\ theorem}\Big)}{=}
\left(\frac{\partial \mu_i}{\partial\theta_i}\right)^{-1} = [v(h(\textbf{z}_i^T\!\mbox{\boldmath $\beta$})) ]^{-1} %\quad \text{(inverse function theorem)}
\end{align*}
\end{itemize}




\end{frame}


\begin{frame}{\S2.2.1 MLE (Score function)}

\begin{itemize}

\item
{\footnotesize Denote $\displaystyle D_i(\mbox{\boldmath $\beta$})=\frac{\partial \mu_i}{\partial\eta_i}
=\frac{\partial h(\eta_i)}{\partial\eta_i}\mid_{\eta_i=\textbf{z}_i^T\!\mbox{\boldmath $\beta$}}$\;
,
$i=1,\cdots, n$ ($1,\cdots, g$ for grouped data).}
\item
{\footnotesize Denote $\displaystyle D(\mbox{\boldmath $\beta$})=\left[\begin{array}{ccc}
\!\!\!D_1(\mbox{\boldmath $\beta$})\!\! & & \\
& \!\!\!\ddots\!\!\! & \\ & & \!\!\!D_n(\mbox{\boldmath $\beta$})\!\!\!\end{array}\right]=\mbox{diag}
\{D_1(\mbox{\boldmath $\beta$}), \cdots, D_n(\mbox{\boldmath $\beta$})\}$}
\item
{\footnotesize Denote $\displaystyle \Sigma(\mbox{\boldmath $\beta$})=\left[\begin{array}{ccc}
\!\!\!\sigma_1^2(\mbox{\boldmath $\beta$})\!\! & & \\
& \!\!\!\ddots\!\!\! & \\ & & \!\!\!\sigma_n^2(\mbox{\boldmath $\beta$})\!\!\!\end{array}\right]=\mbox{diag}
\{\sigma_1^2(\mbox{\boldmath $\beta$}), \cdots, \sigma_n^2(\mbox{\boldmath $\beta$})\}$}


\end{itemize}
In matrix form we get:
\begin{eqnarray*}
%&\!\!\!\stackrel{denoted}{=} \!\!\!& \sum_{i=1}^n \textbf{z}_i D_i(\mbox{\boldmath $\beta$})
%\sigma_i^{-2}(\mbox{\boldmath $\beta$}) \left[y_i-\mu_i(\mbox{\boldmath $\beta$})\right] \\
\textbf{s}(\mbox{\boldmath $\beta$})&\!\!\!=\!\!\!&  \underbrace{\left[\textbf{z}_1, \cdots, \textbf{z}_n\right] }_{Z^T} \left[\!\begin{array}{ccc}
\!\!\!D_1(\mbox{\boldmath $\beta$})\!\! & & \\
& \!\!\!\ddots\!\!\! & \\ & & \!\!\!D_n(\mbox{\boldmath $\beta$})\!\!\!\end{array}\!\right] 
\left[\begin{array}{ccc}\!\!\!\!\sigma_1^2(\mbox{\boldmath $\beta$}) \!\!\! & & \\
&\!\!\!\!\ddots \!\!\!\!& \\ & & \!\!\!\!\sigma_n^2(\mbox{\boldmath $\beta$})\!\!\!\end{array}\right]^{\!\!-1}
\left[\begin{array}{c} \!\!\!y_1\!-\!\mu_1(\mbox{\boldmath $\beta$}) \!\!\!\! \\
\!\!\!\vdots \!\!\!\! \\ \!\!\! y_n\!-\!\mu_n(\mbox{\boldmath $\beta$})\!\!\!\!\end{array}\!\right] \\
&\!\!\!\stackrel{\mbox{}}{=} \!\!\!& Z^TD(\mbox{\boldmath $\beta$})\Sigma^{-1}(\mbox{\boldmath $\beta$})
\left[\textbf{y}-\mbox{\boldmath $\mu$}(\mbox{\boldmath $\beta$})\right] = \sum_{i=1}^n {\bf z}_i D_i(\boldbeta) \sigma_i(\boldbeta)^{-2} (y_i - \mu_i (\boldbeta)).
\end{eqnarray*}
\end{frame}

\begin{frame}{\S2.2.1 MLE (Fisher information)}

The \textbf{Fisher information} matrix is defined as the covariance of the score function at the point $\mbox{\boldmath $\beta$}$: 
\begin{eqnarray*}
F(\mbox{\boldmath $\beta$})&:=& \mbox{cov}\left(\textbf{s}(\mbox{\boldmath $\beta$})\right)= Z^T D \Sigma^{-1}({\boldsymbol \beta}) \Sigma({\boldsymbol \beta})  \Sigma^{-1}({\boldsymbol \beta}) D^T Z  \\
&= &
 Z^T \underbrace{D\Sigma^{-1}({\boldsymbol \beta}) D^T }_{\text{(is diagonal)}}Z = \sum_{i=1}^n \textbf{z}_i D_i^2(\mbox{\boldmath $\beta$})\sigma_i^{-2}(\mbox{\boldmath $\beta$}) 
\textbf{z}_i^T=\sum_{i=1}^n \textbf{z}_i w_i(\mbox{\boldmath $\beta$})\textbf{z}_i^T \\
&=& \left[\textbf{z}_1, \cdots, \textbf{z}_n\right]  \left[\begin{array}{ccc}
\!\!\!w_1(\mbox{\boldmath $\beta$})\!\! & & \\
& \!\!\!\ddots\!\!\! & \\ & & \!\!\!w_n(\mbox{\boldmath $\beta$})\!\!\!\end{array}\right] 
\left[\begin{array}{c} \textbf{z}_1^T \\ \vdots \\ \textbf{z}_n^T\end{array}\right]
=Z^TW(\mbox{\boldmath $\beta$})Z,
\end{eqnarray*}
 where 
 \[
 w_i(\mbox{\boldmath $\beta$}):=D_i^2(\mbox{\boldmath $\beta$})\sigma_i^{-2}(\mbox{\boldmath $\beta$}) \quad \text{ and } \quad 
 W(\mbox{\boldmath $\beta$})=\left[\begin{array}{ccc}
\!\!\!w_1(\mbox{\boldmath $\beta$})\!\! & & \\
& \!\!\!\ddots\!\!\! & \\ & & \!\!\!w_n(\mbox{\boldmath $\beta$})\!\!\!\end{array}\right].
 \]
%  and 
%
%\[ W(\mbox{\boldmath $\beta$})=\left[\begin{array}{ccc}
%\!\!\!w_1(\mbox{\boldmath $\beta$})\!\! & & \\
%& \!\!\!\ddots\!\!\! & \\ & & \!\!\!w_n(\mbox{\boldmath $\beta$})\!\!\!\end{array}\right]
%%=\mbox{diag}
%%\{w_1(\mbox{\boldmath $\beta$}), \cdots, w_n(\mbox{\boldmath $\beta$})\}
%\]

\end{frame}

\begin{frame}{\S2.2.1 MLE (Fisher information in normal LM)}

\begin{wideitemize}
%\item If you have no prior experience with Fisher information, it will become clear that this matrix carries information about the asymptotic covariance structure of our MLE estimate. 


\item Recall in normal linear model (LM), we have 
\[
\hat{\boldsymbol \beta} \sim N({\boldsymbol \beta}, 
\sigma^2({Z}^T{Z} )^{-1} ).
\]

\item $D_i = \frac{\partial \mu_i}{\partial \eta_i} = 1$ implies $W = \sigma^{-2} I_n$;  the scaled \emph{Gram} matrix 
\[
\frac{1}{\sigma^2}{ Z}^T{ Z}
\]
 is in fact the Fisher information for LM. 


\item So  the \emph{inverse} of the Fisher info tells us how well our MLE $\hat{\boldsymbol \beta}$ estimates the true ${\boldsymbol \beta}$ \end{wideitemize}
\end{frame}

\begin{frame}{\S2.2.1 MLE (Observed Fisher information)}
\begin{itemize}
\item A sample-based version of the Fisher info matrix is the {\bf observed information matrix}, defined as
%\item
%The \textbf{observed information} matrix with respect to $\mbox{\boldmath $\beta$}$ is
 {\footnotesize
\begin{eqnarray*}
& \!\!\! F_{\mbox{obs}}(\mbox{\boldmath $\beta$})& \\
& \!\!\!:= \!\!\!& -\frac{\partial^2\ell(\mbox{\boldmath $\beta$})}{
\partial \mbox{\boldmath $\beta$}\partial \mbox{\boldmath $\beta$}^T}=
-\frac{\partial \textbf{s}(\mbox{\boldmath $\beta$})}{\partial \mbox{\boldmath $\beta$}^T} \\
& \!\!\!\underset{\text{(product rule)}}{ =} \!\!\!& -\sum_{i=1}^n \textbf{z}_i \frac{\partial \{D_i(\mbox{\boldmath $\beta$})\sigma_i^{-2}(
\mbox{\boldmath $\beta$})\} }{\partial \mbox{\boldmath $\beta$}^T}[y_i-\mu_i(\mbox{\boldmath $\beta$})]
+\sum_{i=1}^n \textbf{z}_i D_i(\mbox{\boldmath $\beta$})\sigma_i^{-2}(\mbox{\boldmath $\beta$}) \cdot
\underbrace{\frac{\partial\mu_i}{\partial \eta_i}}_{= D_i}\cdot \frac{\partial\eta_i}{\partial \mbox{\boldmath $\beta$}^T} \\
& \!\!\!= \!\!\!& -\sum_{i=1}^n \textbf{z}_i \frac{\partial}{\partial \mbox{\boldmath $\beta$}^T}
\{D_i(\mbox{\boldmath $\beta$})\sigma_i^{-2}(\mbox{\boldmath $\beta$})\}\cdot
\underbrace{[y_i-\mu_i(\mbox{\boldmath $\beta$})] }_{\text{only random part}}+\sum_{i=1}^n \textbf{z}_i D_i^2(\mbox{\boldmath $\beta$})
\sigma_i^{-2}(\mbox{\boldmath $\beta$}) \cdot \textbf{z}_i^T
\end{eqnarray*}}
\item $E[y_i-\mu_i(\mbox{\boldmath $\beta$})] = 0$ implies:
\[E\left[F_{\mbox{obs}}(\mbox{\boldmath $\beta$})\right]=\sum_{i=1}^n \textbf{z}_i 
D_i^2(\mbox{\boldmath $\beta$})\sigma_i^{-2}(\mbox{\boldmath $\beta$}) \textbf{z}_i^T
=F(\mbox{\boldmath $\beta$})=\mbox{cov}\left(\textbf{s}(\mbox{\boldmath $\beta$})\right), \]
i.e.  The expectation of the obs. Fisher info. equals the Fisher info.
\end{itemize}
\end{frame}

\begin{frame}{\S2.2.1 MLE (Natural link)}
\begin{wideitemize}
\item
\underline{When natural link is used},  $\theta_i = \eta_i$ and $h(\cdot) = b'(\cdot)$.

\item $\Rightarrow D_i (\boldbeta) = \frac{\partial h(\eta_i)}{\partial \eta_i}  = b''(\theta_i) = v_i (\mu_i)$.

\item Then,  $\textbf{s}(\mbox{\boldmath $\beta$})$ and $F(\mbox{\boldmath $\beta$})$ further
simplify to
\begin{eqnarray*}
\textbf{s}(\mbox{\boldmath $\beta$})&= & \frac{1}{\phi}\sum_{i=1}^n \omega_i\textbf{z}_i
[y_i-\mu_i(\mbox{\boldmath $\beta$})]=\frac{1}{\phi}Z^T\Omega \left[\textbf{y}-\mbox{\boldmath $\mu$}
(\mbox{\boldmath $\beta$})\right] \\
F(\mbox{\boldmath $\beta$}) &= & \frac{1}{\phi}\sum_{i=1}^n \omega_iv(\mu_i(\mbox{\boldmath $\beta$}))
\textbf{z}_i\textbf{z}_i^T=\frac{1}{\phi}Z^T\Omega V(\mbox{\boldmath $\beta$})Z
\end{eqnarray*}
where 
\begin{equation*}
\Omega \equiv \text{diag}\{\omega_1,\cdots,\omega_n\}  \text{ and }
V({\boldbeta})
=\text{diag}\{v(\mu_1),\cdots, v(\mu_n)\}.
\end{equation*}




\end{wideitemize}
\end{frame}


\begin{frame}{\S2.2.1 MLE (Natural link)}

\begin{itemize}
\item By differentiating the score function one more time wrt to ${\boldsymbol \beta}$ and take a -ve sign, we 
get
\begin{align*}
- \frac{\partial }{ \partial {\boldsymbol \beta}} {\bf s}({\boldsymbol \beta}) &=  \frac{1}{\phi} \sum_{i=1}^n \omega_i {\bf z}_i \frac{\partial \mu_i}{\partial {\boldsymbol \beta}^T} \\
&=  \frac{1}{\phi} \sum_{i=1}^n \omega_i {\bf z}_i  \underbrace{\frac{\partial \mu_i}{\partial \eta_i}}_{v(\mu_i)}  \underbrace{ \frac{\partial \eta_i}{\partial {\boldsymbol \beta}^T} }_{{\bf z}_i^T }  = \frac{1}{\phi} \sum_{i=1}^n \omega_i  v(\mu_i) {\bf z}_i   {\bf z}_i^T.
\end{align*}
So we
have the exact equality
\[
F({\boldsymbol \beta}) = F_{obs}({\boldsymbol \beta}),
\]
not just $E[F_{obs}] = F({\boldsymbol \beta})$. 

\end{itemize}

\end{frame}



\begin{frame}{\S2.2.1 MLE (Fisher scoring)}
\begin{itemize}
\item
MLE $\hat{\mbox{\boldmath $\beta$}}$ is the one maximizing $\ell(\mbox{\boldmath $\beta$})$

\item Often obtained by solving the \textbf{estimating equation} $\textbf{s}(\mbox{\boldmath $\beta$})=0$, using the iterative \textbf{Fisher scoring method}: {\scriptsize
\begin{eqnarray*}
\hat{\mbox{\boldmath $\beta$}}^{(k+1)}&\!\!\!\!\!\!=\!\!\!\!\!\! & \hat{\mbox{\boldmath $\beta$}}^{(k)}+
F^{-1}(\hat{\mbox{\boldmath $\beta$}}^{(k)})\cdot \textbf{s}(\hat{\mbox{\boldmath $\beta$}}^{(k)}) \\
&\!\!\!\!\!\!=\!\!\!\!\!\! & \hat{\mbox{\boldmath $\beta$}}^{(k)}+ 
\left[Z^TW(\hat{\mbox{\boldmath $\beta$}}^{(k)})Z\right]^{-1}
 Z^TD(\hat{\mbox{\boldmath $\beta$}}^{(k)})\Sigma^{-1}(\hat{\mbox{\boldmath $\beta$}}^{(k)})
\left[\textbf{y}-\mbox{\boldmath $\mu$}(\hat{\mbox{\boldmath $\beta$}}^{(k)})\right] \\
&\!\!\!\!\!\!=\!\!\!\!\!\! & \hat{\mbox{\boldmath $\beta$}}^{(k)}+ 
\left[Z^TW(\hat{\mbox{\boldmath $\beta$}}^{(k)})Z\right]^{-1}
 Z^TD^2(\hat{\mbox{\boldmath $\beta$}}^{(k)})\Sigma^{-1}(\hat{\mbox{\boldmath $\beta$}}^{(k)})
D^{-1}(\hat{\mbox{\boldmath $\beta$}}^{(k)}) 
\left[\textbf{y}-\mbox{\boldmath $\mu$}(\hat{\mbox{\boldmath $\beta$}}^{(k)})\right] \\
&\!\!\!\!\!\!=\!\!\!\!\!\! & \left[Z^TW(\hat{\mbox{\boldmath $\beta$}}^{(k)})Z\right]^{-1}\left\{
Z^TW(\hat{\mbox{\boldmath $\beta$}}^{(k)})Z \hat{\mbox{\boldmath $\beta$}}^{(k)}
+Z^TW(\hat{\mbox{\boldmath $\beta$}}^{(k)}) D^{-1}(\hat{\mbox{\boldmath $\beta$}}^{(k)})
\left[\textbf{y}-\mbox{\boldmath $\mu$}(\hat{\mbox{\boldmath $\beta$}}^{(k)})\right]\!\right\} \\
&\!\!\!\!\!\!=\!\!\!\!\!\!& \left[Z^TW(\hat{\mbox{\boldmath $\beta$}}^{(k)})Z\right]^{-1}\left\{
Z^TW(\hat{\mbox{\boldmath $\beta$}}^{(k)})\left[Z \hat{\mbox{\boldmath $\beta$}}^{(k)}
+ D^{-1}(\hat{\mbox{\boldmath $\beta$}}^{(k)})
\left[\textbf{y}-\mbox{\boldmath $\mu$}(\hat{\mbox{\boldmath $\beta$}}^{(k)})\right]\right]\right\} \\
&\!\!\!\!\!\!\stackrel{denoted}{=}\!\!\!\!\!\! & \left[Z^TW(\hat{\mbox{\boldmath $\beta$}}^{(k)})Z\right]^{-1}
Z^TW(\hat{\mbox{\boldmath $\beta$}}^{(k)})\cdot \tilde{\textbf{y}}^{(k)}
\end{eqnarray*}}
where $\tilde{\textbf{y}}^{(k)}=(\tilde{y}_1^{(k)},\cdots, \tilde{y}_n^{(k)})^T
=Z \hat{\mbox{\boldmath $\beta$}}^{(k)}+ D^{-1}(\hat{\mbox{\boldmath $\beta$}}^{(k)})
\left[\textbf{y}-\mbox{\boldmath $\mu$}(\hat{\mbox{\boldmath $\beta$}}^{(k)})\right]$
is the \textbf{working observation vector} at iteration $k$, $k=0,1,2,\cdots$.
\end{itemize}
\end{frame}

\begin{frame}{\S2.2.1 MLE (Fisher scoring)}

\begin{wideitemize}
\item If  we use $F_{obs} (\cdot)$ instead of $F (\cdot)$, we get the standard Newton-Raphson procedure; the two procedures coincide if the natural link  is used.

\item For non-canonical links, $F(\boldbeta)$ is generally easier to evaluate than $F_{obs} (\boldbeta)$ (the latter involves taking second derivatives) and  positive definite enough over a large region in the parameter space (easy to invert numerically).
\end{wideitemize}


\end{frame}
\begin{frame}{\S2.2.1 MLE (Fisher scoring) }
\begin{enumerate}
\item
Start with an initial estimate $\hat{\mbox{\boldmath $\beta$}}^{(0)}$, often chosen as the unweighted
least squares (LS) estimate from $(g(y_i),\textbf{z}_i)$, $i=1,\cdots,n$, where $g(\cdot)$ is the link (modified if $g(y_i)$ is undefined, such as when $y_i = 0$ in Poisson regression, where the natural link is $\log(\cdot)$).

\ \
\item
For $k=0,1,2,\cdots$ compute 
$$\hat{\mbox{\boldmath $\beta$}}^{(k+1)}=\left[Z^TW(\hat{\mbox{\boldmath $\beta$}}^{(k)})Z\right]^{-1}
Z^TW(\hat{\mbox{\boldmath $\beta$}}^{(k)})\cdot \tilde{\textbf{y}}^{(k)}$$
until $\frac{\|\hat{\boldsymbol \beta}^{(k+1 )} - \hat{\boldsymbol \beta}^{(k)}\|}{\|\hat{\boldsymbol \beta}^{(k)}\|} \leq \epsilon $ for some prechosen small $\epsilon > 0$ (``convergence").

\ \

\item Also called a  \textbf{iteratively reweighted least squares (IRLS)} method. 
\end{enumerate}

\end{frame}


\begin{frame}{\S2.2.1 MLE (Fisher scoring) }

\begin{itemize}

\item 
Each component of the working observation vector is a linearization of $g(y_i)$ at $\boldbeta = \hat{\boldbeta}^{(k)}$:

\begin{align*}
g(y_i) &\approx g(\mu) +g'(\mu) (y_i - \mu) \\
&= {\bf z}_i^T{\boldsymbol \beta} +  \left( \underbrace{\frac{\partial h}{ \partial \eta} ({\bf z}_i^T{\boldsymbol \beta} )}_{= D_i ({\boldsymbol \beta})}\right)^{-1}  (y_i - \mu_i ({\boldsymbol \beta}))
\end{align*}
(Generalized Linear Models, 2nd ed, McCullagh and Nelder, p.40)


\end{itemize}
\end{frame}

\begin{frame}{\S2.2.1 MLE (Asymptotic properties)}
Under  ``regularity" assumptions, we have these nice properties:
\begin{wideitemize}
\item
\textit{\textcolor{red}{Asymptotic existence and uniqueness}}: With probability tending to $1$ as $n \rightarrow \infty$,  the MLE $\hat{\mbox{\boldmath $\beta$}}$ exists
and is (locally) unique. 
\item
\textit{\textcolor{red}{Consistency}}. If $\mbox{\boldmath $\beta$}$ is the ``true'' value, then 
$\hat{\mbox{\boldmath $\beta$}}\rightarrow \mbox{\boldmath $\beta$}$ in probability (\textit{weak consistency}),
or $\hat{\mbox{\boldmath $\beta$}}\rightarrow \mbox{\boldmath $\beta$}$ with probability 1 (
\textit{strong consistency}).

\item
\textit{\textcolor{red}{Asymptotic normality}}. $\hat{\mbox{\boldmath $\beta$}}\stackrel{a}{\sim} N
\left(\mbox{\boldmath $\beta$}, F^{-1}(\boldbeta)\right)$, i.e. $\hat{\boldbeta}$ is approximately normal with mean $\boldbeta$ and  covariance matrix $F^{-1} (\boldbeta)$. Replacing $\boldbeta$ by $\hat{\boldbeta}$ gives an approximate covariance that can be used in practice, i.e. 
\[
\text{cov}(\hat{\boldbeta}) \stackrel{a}{=} F^{-1}(\hat{\boldbeta}).
\]


\end{wideitemize}

\begin{enumerate}
\item $\stackrel{a}{=}$  means ``approximately equal to" 
\item $\stackrel{a}{\sim}$ means ``approximately distributed as"
\end{enumerate}
\end{frame}




% new frame

\begin{frame} {\S2.2.1  MLE (Asymptotic properties)}

The following discussion comes from F \&T (p.46 and  appendix A.2), but cannot replace a rigorous course on asymptotic statistics \cite{ferguson1996}.

\ \

\begin{itemize}

\item
Note that both the score function and Fisher information matrix are sums of $n$ independent terms:
\begin{align*}
{\bf s}({\boldsymbol \beta}) &= {\bf s}_n ({\boldsymbol \beta}) = \sum_{i=1}^n {\bf z}_i^T D_i(\boldbeta) \sigma_i(\boldbeta)^{-2} (y_i - \mu_i) \\
F({\boldsymbol \beta}) &= F_n({\boldsymbol \beta})  =  \sum_{i=1}^n {\bf z}_i w_i(\boldbeta){\bf z}_i^T \\
\end{align*}
\item Caveat: The summands are generally not identically distributed because ${\bf x}_i$ are different for different $i$.

\end{itemize}

\end{frame}

% new frame
\begin{frame} {\S2.2.1  Consistency}
\begin{itemize}
\item Since $E[y_i] = \mu_i(\boldbeta)$ for the ``true" $\boldbeta$, it follows that
\[
E[{\bf s}({\boldsymbol \beta})] = \sum_{i=1}^n {\bf z}_i^T D_i(\boldbeta) \sigma_i(\boldbeta)^{-2}  E[y_i - \mu_i] = 0
\]


\item By some law of large numbers, ${\bf s} (\boldbeta)/n =  {\bf s}_n (\boldbeta)/n \longrightarrow 0$ in probability, at the true $\boldbeta$. Since 
\[
{\bf s} (\hat{\boldbeta}) = 0
\]
 for the MLE $\hat{\boldbeta}$, it can be expected that 
 \[
 \hat{\boldbeta} \rightarrow \boldbeta \text{ in probability,}
 \]
   i.e.  we have the  consistency as claimed. 
\end{itemize}
\end{frame}

\begin{frame} {\S2.2.1  Asymptotic normality}


\begin{itemize}



\item One ``standard" regularity assumption is that
\[
%\frac{F_n(\boldbeta)}{n} \text{ converges to a positive definite limit, i.e.  }
 \frac{F(\boldbeta)}{n}  \longrightarrow F_{limit} \text{ for some limiting positive definite matrix } F_{limit}.
\]

\item Along with other weak conditions, the score function is asymptotically normal:
\[
n^{-1/2} {\bf s} (\boldbeta) \longrightarrow_d N(0, F_{limit}(\boldbeta))
\]


\item By Taylor expansion, we have 
\begin{equation*}
0 = {\bf s} (\hat{\boldbeta}) \approx  {\bf s}(\boldbeta) -  F_{obs} (\boldbeta)(\hat{\boldbeta} - \boldbeta) \approx   {\bf s}(\boldbeta) -  F (\boldbeta)(\hat{\boldbeta} - \boldbeta).
\end{equation*}
Therefore, one can expect
\[
\sqrt{n} (\hat{\boldbeta} - \boldbeta) \approx \sqrt{n} F^{-1}(\boldbeta){\bf s} (\boldbeta) \rightarrow_d N(0, F_{limit}^{-1} (\boldbeta)),
\]
and hence the approximate normality $\hat{\boldbeta} \stackrel{a}{\sim} N
\left(\mbox{\boldmath $\beta$}, F^{-1}(\hat{\mbox{\boldmath $\beta$}})\right)$.
\end{itemize}



\end{frame}

%new frame
\begin{frame}{\S2.2.1 Asymptotic efficiency of MLE}

\begin{itemize}
\item In mathematical statistics, one fundamental result is the \textcolor{red}{Cramer-Rao lower bound}.

\item It says that, under suitable regularity conditions, for any unbiased estimator  $\bolddelta (y_1, \dots, y_n; {\bf x}_1, \dots, {\bf x}_n)$ of $\boldbeta$, it must be the case that
\[
\text{cov}(\bolddelta) - F^{-1} (\boldbeta) \geq 0,
\]
i.e., $\text{cov}(\bolddelta) - F^{-1} (\boldbeta)$, as a matrix, is positive semi-definite.

\item In particular, for any vector ${\bf c} \in \mathbb{R}^p$, it is true that
\[
\text{Var} ({\bf c}^T \bolddelta) = {\bf c}^T \text{cov}(\bolddelta)  {\bf c} \geq  {\bf c}^T  F^{-1} (\boldbeta){\bf c}.
\]
i.e. $F^{-1} (\boldbeta)$ is the ``lowest" variance an unbiased est'r of $\boldbeta$ can achieve. 
\item That $\hat{\boldbeta} \stackrel{a}{\sim} N
\left(\mbox{\boldmath $\beta$}, F^{-1}(\boldbeta)\right)$ means 
\[
\text{ the MLE is \emph{asymptotically unbiased and efficient}.}
\]
\end{itemize}





\end{frame}


\begin{frame} {\S2.2.1  Delta method}

\begin{itemize}

 
\item Let  $t: \mathbb{R}^p \rightarrow \bR$ be any smooth (possibly multivariate) function in  $\boldbeta$.

\item The \emph{delta method} says, if 
\[
\sqrt{n}(\hat{\boldbeta} - \boldbeta) \longrightarrow_d  N(0, F_{limit}^{-1} (\boldbeta)),
\]
then 
\[
\sqrt{n}(t(\hat{\boldbeta}) - t(\boldbeta)) \longrightarrow_d  N( 0,   \nabla t(\boldbeta)^T F_{limit}^{-1} (\boldbeta)  \nabla t(\boldbeta)),
\]
where $\nabla t(\cdot)$ here denotes the gradient of $t(\cdot)$.

\item In our context, it means 
\[
t(\hat{\boldbeta}) \sim_a N(t(\boldbeta) ,   \nabla t(\boldbeta)^T F^{-1}(\boldbeta) \nabla t(\boldbeta)).
\]




\item One can then construct confidence interval for $t(\boldbeta)$ based on the estimate $t(\hat{\boldbeta})$, using the estimated  variance $\nabla t(\hat{\boldbeta})^T F^{-1}(\hat{\boldbeta}) \nabla t(\hat{\boldbeta})$. 

\end{itemize}


\end{frame}

\begin{frame} {\S2.2.1 For grouped data..}


\begin{wideitemize}





\item Another regularity assumption concerns grouped data: Assume a fixed number $g$  of groups and that for each $i \in \{1, \dots, g\}$, $\frac{n_i}{n} \rightarrow \lambda_i $, where $\lambda_i >0$. 

\item However, this assumption is likely violated if there is a continuous covariate. 



\end{wideitemize}



\end{frame}



\begin{frame} {\S2.2.1 Estimation of $\phi$}
\begin{wideitemize}
\item We also note that $\phi$ only appears as a factor in $W$, and falls out from the expression

 \[\hat{\mbox{\boldmath $\beta$}}^{(k+1)}=\left[Z^TW(\hat{\mbox{\boldmath $\beta$}}^{(k)})Z\right]^{-1}
Z^TW(\hat{\mbox{\boldmath $\beta$}}^{(k)})\cdot \tilde{\textbf{y}}^{(k)}\]
\item But we need $\phi$ for computing the Fisher info, to construct the approximate covariance of the MLE. 

\item\[\hat{\phi}=\frac{1}{g-p}\sum_{i=1}^g\frac{(\bar{y}_i-\hat{\mu}_i)^2}{v(\hat{\mu}_i)/n_i}\]
is a consistent estimator of $\phi$, with $\hat{\mu}_i = h({\bf z}_i^T \hat{\boldbeta})$.
\end{wideitemize}
\end{frame}



\begin{frame} {\S2.2.1 Estimation of $\phi$}
\begin{itemize}


\item A conventional  (no grouping) moment estimator for $\hat{\phi}$ is
\[
\frac{1}{n - p} \sum_{i = 1}^n (y_i - \hat{\mu}_i)^2 /v_i(\hat{\mu}_i)
 \]
(MuCullagh and Nelder 1989, p.328).  \end{itemize}
\end{frame}

\subsection{\S2.2.2 Linear hypothesis testing in GLM}

\begin{frame}{\S2.2.2 Testing linear hypotheses}
\begin{itemize}
\item
Most testing problems concerning $\mbox{\boldmath $\beta$}$ have the form of \textit{linear hypothesis}:
$$H_0: C\mbox{\boldmath $\beta$}=\mbox{\boldmath $\xi$} \quad\mbox{vs.}\quad
H_1: C\mbox{\boldmath $\beta$}\neq\mbox{\boldmath $\xi$}$$
where $C$ is a known matrix of full row rank $s\leq p$, and $\mbox{\boldmath $\xi$}$ is a known vector.
\item
A special case is 
$$H_0: \mbox{\boldmath $\beta$}_r=0 \quad\mbox{vs.}\quad
H_1: \mbox{\boldmath $\beta$}_r\neq 0$$
where $\mbox{\boldmath $\beta$}_r$ is a sub-vector of $\mbox{\boldmath $\beta$}$. This compares the full GLM  with its sub-model that assumes $\boldbeta_r = 0$.
\end{itemize}
\end{frame}


\begin{frame}{\S2.2.2 Likelihood ratio (LR), Wald and score tests (1)}
\begin{wideitemize}
\item
Suppose we are to test the linear hypothesis:
\[
H_0: C\mbox{\boldmath $\beta$}=\mbox{\boldmath $\xi$}
\;\mbox{vs }\; H_1: C\mbox{\boldmath $\beta$}\neq\mbox{\boldmath $\xi$}
\]

\item Let $\tilde{\boldbeta} \equiv \underset{\boldbeta: C \boldbeta = {\boldsymbol \xi}}{\text{argmax}} \hspace{1mm}\ell(\boldbeta)$ be the (restricted) MLE of
$\boldbeta$  under $H_0$

\item Let  $\hat{\boldbeta} \equiv \underset{\boldbeta}{\text{argmax}} \hspace{1mm} \ell(\boldbeta)$ be the (unrestricted) MLE of
$\boldbeta$ 


\end{wideitemize}
\end{frame}

\begin{frame}{\S2.2.2 Likelihood ratio (LR), Wald and score tests (2)}
\begin{enumerate}
\item
\textbf{LR test statistic}:
 \[\lambda=-2\left\{\ell(\tilde{\mbox{\boldmath $\beta$}}) -\ell(\hat{
\mbox{\boldmath $\beta$}})\right\}.
\]
 
\ \
 
%
%$\lambda$ needs to be divided by $\hat{\phi}$ for over-dispersion models.
\item
\textbf{Wald test statistic}: \[W=\left(C\hat{\mbox{\boldmath $\beta$}}-\mbox{\boldmath $\xi$}\right)^T
\left[CF^{-1}(\hat{\mbox{\boldmath $\beta$}})C^T \right]^{-1}
\left(C\hat{\mbox{\boldmath $\beta$}}-\mbox{\boldmath $\xi$}\right)
\] 

\ \

\item
\textbf{Score test statistic}: $U=\textbf{s}^T(\tilde{\mbox{\boldmath $\beta$}})F^{-1}(\tilde{\mbox{\boldmath $\beta$}})
\textbf{s}(\tilde{\mbox{\boldmath $\beta$}})$.
\end{enumerate}

\end{frame}




\begin{frame}{\S2.2.2 Likelihood ratio (LR), Wald and score tests (3)}
\begin{wideitemize}
%\item  The  exponential model is assumed for the LR test; not defined for over-dispersion model such as the quasi likelihood (to be covered). 
%
%\item 
%Wald and Score test statistics are properly defined for over-dispersion models, since they are based on the score ${\bf s}$ (Again, more on this later)


\item Wald and Score are computationally attractive; both eliminate the need to compute one of $\tilde{\boldsymbol \beta}$ or $\hat{\boldsymbol \beta}$ 

\begin{wideitemize}

\item
 Wald is good when $H_1$ has been fitted.
 
 \item Score is good when $H_0$ has been fitted. 
 
 \end{wideitemize}
\item When the null $H_0$ is true, under  suitable ``regularity conditions", 

\[
\lambda, W \text{ and }
U \text{ all follow a } \chi^2(s) \text{ distribution asymptotically.}
\]

\end{wideitemize}

\end{frame}


\begin{frame} {Some theoretical backdrop for LR, Wald and score}
\begin{wideitemize}
\item LR, Wald and score are asymptotically equivalent, i.e. they are close to each other as $n \rightarrow \infty$. 

\item Hence they  have the same limiting null distribution.

\item For  a simple hypothesis $H_0: \boldbeta = \boldbeta_0$, upon Taylor expansion, we have
\begin{align*}
\lambda &= - 2(\ell(\boldbeta_0) - \ell(\hat{\boldbeta})) \\
&=  \underbrace{(\hat{\boldbeta} - \boldbeta_0)^T F(\hat{\boldbeta})(\hat{\boldbeta} - \boldbeta_0)}_{W}+ O_p(n^{-1/2}) \\
& =  \underbrace{{\bf s}^T(\boldbeta) F^{-1}(\boldbeta) {\bf s}(\boldbeta) }_{U} + O_p(n^{-1/2})
\end{align*}
 (F\&T, p.440).

\end{wideitemize}



\end{frame}

\begin{frame}{Some theoretical backdrop for LR, Wald and score}
\begin{wideitemize}

\item Discrepancies between the three statistics in finite sample $n$:
\begin{figure}
 \centering
\includegraphics[width=0.5\textwidth]{quadratic}
% \label{figure:example}
\end{figure}

\item For composite null hypothesis, $\lambda = -2\left\{\ell(\tilde{\mbox{\boldmath $\beta$}}) -\ell(\hat{
\mbox{\boldmath $\beta$}})\right\}$ can be expressed as the difference of two quadratic forms, which is asymptotically $\chi^2$ with the correct degree of freedom \cite{CH1974}. 

\end{wideitemize}
\end{frame}

%
%\begin{frame}{Some theoretical backdrop for LR, Wald and score}
%\begin{wideitemize}
%
%\item Discrepancies between the three statistics in finite sample $n$:
%\begin{figure}
% \centering
%\includegraphics[width=0.5\textwidth]{quadratic}
%% \label{figure:example}
%\end{figure}
%
%\item For composite null hypothesis, $\lambda = $
%
%\end{wideitemize}
%\end{frame}

\subsection{\S2.2.3 Deviance, residuals and goodness-of-fit measure in GLM} \label{sec:2.2.3}

\begin{frame} {\S2.2.3 Model adequacy measure in GLM}
\begin{wideitemize}
\item One can assess the model adequacy, or \emph{goodness-of-fit},  of a GLM. 

\item
In the normal linear model (LM), goodness-of-fit  is measured by the sum of squared residuals 
$\mbox{SSE}\!=\!\sum_{i=1}^n(y_i-\hat{\mu}_i)^2$. 

\item SSE puts equal weights on the squared residuals $(y_i-\hat{\mu}_i)^2$

\item $\Rightarrow$ not quite reasonable as a goodness-of-fit measure if $\mbox{var}(y)$ 
varies with $\mu$, such as in GLM.

\item Two goodness-of-fit measures for GLMs: {\bf deviance} and {\bf Pearson statistics}.
\end{wideitemize}
\end{frame}


\begin{frame}{\S2.2.3 Deviance }
\begin{itemize}

\item
The {\bf scaled deviance} $D^*$ is twice the difference between the maximum log-likelihoods achievable by the \underline{saturated model} and by
the \underline{GLM}, i.e. for $\boldtheta = (\theta_1, \dots, \theta_g)$,
\[
D^*=2\left[\ell(\tilde{\boldtheta},\phi ; \bar{\bf y})
-\ell(\mbox{\boldmath $\hat{\theta}$},\phi ; \bar{\bf y})\right],
\]
where
\[
\ell({\boldsymbol \theta}, \phi; \bar{\bf y}) = \sum_{i =1}^g \log f( \bar{y}_i |\theta_i, \phi, \omega_i)
\]
for the grouped data vector $\bar{\bf y} = (\bar{y}_1, \dots, \bar{y}_g)$ and the exponential density form $f( \cdot |\theta_i, \phi, \omega_i)$ of the GLM.



%\item (Remark: we have grouped data in mind here)


\end{itemize}
\end{frame}

\begin{frame}{\S2.2.3 Deviance }

\begin{itemize}
\item

 $\mbox{\boldmath $\hat{\theta}$}$ is the MLE of
$\mbox{\boldmath $\theta$}$ under the GLM, i.e., for  $\hat{\mu}_i = h({\bf z}_i^T \hat{\boldbeta})$,
\[
\hat{\boldtheta} = (b'^{-1} (\hat{\mu}_1), \dots, b'^{-1} (\hat{\mu}_g)),
\]
where $b'^{-1}(\cdot)$ means the inverse function of $b'(\cdot)$. In particular, we have the \emph{structural assumptions} $\mu_i = h({\bf z}_i^T \boldbeta)$ in mind when we refer to the ``GLM".


\item $\tilde{\boldtheta}=(\tilde{\theta}_1,\cdots, \tilde{\theta}_g)$ is defined as
\[
(\tilde{\theta}_1,\cdots, \tilde{\theta}_g) \equiv 
\left(b'^{-1}(\bar{y}_1),\cdots,b'^{-1}(\bar{y}_g) \right);
\]
it is the MLE in the {\bf saturated model}.

\item The {\bf saturated model}: Defined by the exponential family density of the GLM, but \textcolor{red}{without} the structural assumptions. In particular, it is parametrized by $\boldtheta$ directly, not by $\boldbeta$.

\end{itemize}

\end{frame}

\begin{frame}{\S2.2.3 Deviance  }
\begin{itemize}

\item
By differentiating $\ell$ with respect to $\boldtheta = (\theta_1, \dots, \theta_g)$ and equating  to zero, we get
\[
\frac{\partial \ell}{ \partial \theta_i} =  \frac{\bar{y}_i - b'(\theta_i)}{ \phi} \omega_i = 0 \quad \text{ for each } \quad i = 1, \dots, g;
\]
so the MLE $\tilde{\boldtheta}$ of the saturated model
must satisfy $\bar{y}_i=b^\prime(\tilde{\theta}_i)$. 


\item

The maximum likelihood of the saturated model must be no less than that of the  GLM with structural assumptions, because the former maximizes the likelihood in the least restrictive way.
Therefore, 

\[
D^*\geq 0.
\] 
\item

By convention
\[
D:=\phi D^*
\]
is defined as  the \textbf{deviance}.



\end{itemize}
\end{frame}




\begin{frame}{\S2.2.3 Deviances for several important distributions}
\textbf{Deviances for several important distributions in GLM.}
\begin{eqnarray*}
\mbox{Normal}\!\!\! & N(\mu,\sigma^2): & D=
\sum_{i=1}^n (y_i-\hat{\mu}_i)^2 \text{ (coincides with SSE)}\\ 
\mbox{Poisson}\!\!\! &\mbox{Poi}(\mu): &  D=2\sum_{i=1}^n
\left\{y_i\ln\frac{y_i}{\hat{\mu}_i}-(y_i-\hat{\mu}_i)\right\} \\
\mbox{binomial}\!\!\! & \mbox{Bin}(m,\mu): & D=2\sum_{i=1}^n
\left\{y_i\ln\frac{y_i}{\hat{\mu}_i}\!+\!(\!m_i\!-\!y_i\!)\ln\frac{m_i\!-\!y_i}{m_i\!-\!\hat{\mu}_i}\right\} \\
\mbox{gamma} \!\!\!& \mbox{Gam}(\mu,\nu): & D=2\sum_{i=1}^n
\left\{-\ln\frac{y_i}{\hat{\mu}_i}+\frac{y_i-\hat{\mu}_i}{\hat{\mu}_i}\right\}
\end{eqnarray*}
\end{frame}



\begin{frame}{\S2.2.3 Pearson statistics }
\begin{wideitemize}

\item
The \textbf{Pearson $\chi^2$ statistic} takes the form 
\[
\chi^2=\sum_{i=1}^g \frac{(\bar{y}_i-\hat{\mu}_i)^2}{v(\hat{\mu}_i)/n_i}
\]
where  $n_i$ is the size of group $i$.

\item An intuitive generalization of the SSE

\item This can be derived as the {\bf score test} for comparing the GLM with the saturated model \cite{smyth2003pearson, Lovison2005}.

\item  So the relationship between the Pearson and deviance  statistics is the relationship between the score test and likelihood ratio statistics. 
\end{wideitemize}
\end{frame}



\begin{frame}{\S2.2.3  Asymptotics for deviance and Pearson statistics  }
\begin{wideitemize}

\item
Asymptotically, 
\[
\phi^{-1}\chi^2, \phi^{-1}D\sim_a\chi^2(g-p),
\] 

where $p$ is the number of parameters in $\boldbeta$ 
when
\begin{itemize}
\item  the GLM  is indeed true,
\item the group sizes $n_i$ are all large enough, and 
\item under other ``usual" regularity conditions.
\end{itemize}


\item In the canonical normal linear model 
\[
y_i  = {\bf z}_i^T \boldbeta + \epsilon_i
\]
for i.i.d. $\epsilon_1, \dots, \epsilon_n \sim N(0, \sigma^2)$,
\[
\phi^{-1}\chi^2= \phi^{-1}D = \frac{ \sum_{i=1}^n(y_i-{\bf z}_i^T \hat{\boldbeta})^2}{\sigma^2} \sim \chi^2 (n - p),
\]
where the last distributive law is exact  \cite[Sec 4]{SeberLee2003}.

\end{wideitemize}
\end{frame}


\begin{frame}{Digression: Wilks' theorem (1) }
\begin{wideitemize}

\item
The chi-square asymptotics above are based on what is known as {\bf Wilks' theorem}.


\item In a general setup with log-likelihoods $\ell(\theta)$ indexed by a parameter $\theta$ in the space $\Theta$, suppose we test
\[
H_0: \theta \in \Theta_0 ,
\]
 where $\Theta_0$ is a lower dimensional subspace of the ambient space $\Theta$.





\end{wideitemize}
\end{frame}

\begin{frame}{Digression: Wilks' theorem (2) }
\begin{wideitemize}

\item Let 
\[
\lambda = \sup_{\theta \in \Theta} \sum_i \ell_i(\theta) -  \sup_{\theta \in \Theta_0} \sum_i \ell_i(\theta)
\] be the log maximized likelihood ratio, where $\ell_i$ is the log-likelihood for an individual sample.

\item Under suitable regularity conditions and when $H_0$ is true, Wilks' theorem says  
\[
2 \lambda \text{ converges in distribution to }\chi^2(d-d_0)
\]
 as samples grow, where $d = \dim (\Theta)$ and $d_0 = \dim (\Theta_0)$ are the respective dimensions of the two spaces.


\item For the goodness-of-fit tests above, $d = g$ and $d_0 = p$.




\end{wideitemize}
\end{frame}


\begin{frame}{\S2.2.3  Asymptotics for deviance and Pearson statistics  }
\begin{wideitemize}



\item
If $n$ is large but $n_i$'s are small (in particular when $n_i=1$, as in ungrouped data), the chi-square asymptotics 
\[
\phi^{-1}\chi^2, \phi^{-1}D\sim_a\chi^2(g-p),
\]
 may not be valid. 

\item In particular, large $D$ or $\chi^2$ may not necessarily invalidate the GLM.  


\item If $n_i$ are too small, the dimension of the ambient parameter space (for the saturated model)  keeps growing. For the null model (the GLM), the parameter space $\Theta_0 = \Theta(\boldbeta)$ is a (one-to-one, at least locally) function of $\boldbeta$, so it must have dimension $p$.


\end{wideitemize}
\end{frame}



\begin{frame}{\S2.2.3  Asymptotics for deviance and Pearson statistics  }
\begin{wideitemize}



\item $\Rightarrow$ $\phi^{-1}D$ (and asymp. equivalently, $\phi^{-1}\chi^2$) cannot possibly converge to a stable chi-square distribution even if the GLM is true!


\item McCullagh and Nelder (p.118-119) also discussed how the asymptotic theory breaks down


\item McCullagh and Nelder (p.34) also makes the general remark:
 \ \

``\emph{The deviance has a general advantage as a measure of discrepancy in that it is additive for nested sets of models if MLE's are used, where as $\chi^2$ in general is not. However, $\chi^2$ may sometimes be preferred because of its more direct interpretation.} "
\end{wideitemize}
\end{frame}





\begin{frame}{\S2.2.3  Asymptotics for testing linear hypotheses under the GLM, revisited }

\begin{wideitemize}
\item Wilks' theorem can also be used to understand the  chi-square asymptotic for testing 
\[
H_0: C{\boldmath \boldbeta}={\boldmath \xi} 
\] 
where $C$ is a $s \times p$ matrix with full row rank $s$ (so $s \leq p$), against the full GLM model with no restriction on $\boldbeta$.

\item Recall that 
$\lambda, W$ and
$U$  follow the $\chi^2(s)$ distribution asymptotically under the null. How do we get the degree $s$ in the limiting $\chi^2$?

\item The unrestricted full GLM model has dimension $p$.
\end{wideitemize}


\end{frame}


\begin{frame}{\S2.2.3  Asymptotics for testing linear hypotheses under the GLM, revisited }

\begin{itemize}


\item The restricted GLM under $H_0$ has the same dimension as the space of all $\boldbeta$ satisfying $C{\boldmath \boldbeta}={\boldmath \xi}$.

\item Since $C$ is a $s \times p$ matrix with full row rank, by elementary linear algebra,  the $s \times s$ matrix $C C^T$ is invertible. Hence, $H_0$ is  equivalent to
\[
C{\boldmath \boldbeta} -  C C^T (C C^T)^{-1} \xi = C (\boldbeta - \xi') = 0,
\]
where $\xi' := C^T (C C^T)^{-1} \xi$.

\item Hence, under $H_0$, $\boldbeta - \xi'$ lies in the linear null space of $C$,which has dimension $p-s$.

\item So $\boldbeta$ lies in this $p-s$ dimensional space shifted by $\xi'$, which has the same dimension.


\item The chi-square df is hence $p - (p -s) = s$, by Wilk's theorem.
\end{itemize}


\end{frame}


\begin{frame}[allowframebreaks]
  \frametitle{References}    
  \begin{thebibliography}{10}    

%   \beamertemplatearticlebibitems
%  \bibitem[Lang, 1990]{Lang2014}
%    Lang, Joseph B
%    \newblock The Pearson Score Statistic for Multinomial- Poisson Models
%    \newblock \emph{Communications in Statistics-Theory and Methods 43.21 (2014): 4471-4491}
  \beamertemplatearticlebibitems
  \bibitem[Smyth, 2003]{smyth2003pearson}
    Smyth, Gordon K
    \newblock Pearson's goodness of fit statistic as a score test statistic
    \newblock {\em Lecture notes-monograph series (2003): 115-126.}
    
      \beamertemplatearticlebibitems
  \bibitem[Lovison, 2005]{Lovison2005}
   Lovison, Gianfranco. 
    \newblock On Rao score and Pearson X 2 statistics in generalized linear models.
    \newblock {\em Statistical Papers 46 (2005): 555-574.}
    
      \beamertemplatearticlebibitems
  \bibitem[Cox and Hinkley 1974]{CH1974}
    David R. Cox, David V. Hinkley
    \newblock Theoretical Statistics
    \newblock {\em Chapman \& Hall (1974).}
    
       \beamertemplatearticlebibitems
  \bibitem[Ferguson 1996]{ferguson1996}
    Thomas S. Ferguson
    \newblock A Course in Large Sample Theory
    \newblock {\em Chapman \& Hall (1996).}
    
         \beamertemplatearticlebibitems
  \bibitem[Seber and Lee 2003]{SeberLee2003}
    George A. F. Seber \& Alan J. Lee
    \newblock Linear Regression Analysis, Second Edition
    \newblock {\em Wiley (2003).}

%    
%      \beamertemplatearticlebibitems
%  \bibitem[Moore, 1977]{moore}
%    Moore David S
%    \newblock Generalized inverses, Wald's method, and the construction of chi-squared tests of fit
%    \newblock {\emph Journal of the American Statistical Association 72.357 (1977): 131-137}
%    
    
%            \beamertemplatearticlebibitems
%  \bibitem[Seber and Lee, 2003]{seberleelinear}
%   Seber, G. A., \& Lee, A. J.
%    \newblock Linear regression analysis, Second edition.
%    \newblock {\emph John Wiley \& Sons. 2003. }
%    
%    
%            \beamertemplatearticlebibitems
%  \bibitem[Rao and Mitra, 1972]{radhakrishna1972generalized}
% Radhakrishna Rao, C., and Sujit Kumar Mitra
%    \newblock Generalized inverse of a matrix and its applications
%    \newblock {In Proceedings of the Berkeley Symposium on Mathematical Statistics and Probability, vol. 1, pp. 601-620. University of California Press, Berkeley, 1972. }
    
    
  \end{thebibliography}
\end{frame}
\end{document}

